\section{结果与讨论}\label{sec:results}

利用上一节提取的RITD，我们考察$\pi$介子质量与晶格间距依赖。
图~\ref{fig:RITD-ensembles}上图给出 a12m220、a12m310 与 a15m310 三个系综在提升动量约2~GeV处$\eta_s$的RITD随 Wilson 线长度$z$的变化。
我们未看到明显的晶格间距依赖。
图~\ref{fig:RITD-ensembles}下图给出同样三个系综在提升动量约1.3~GeV处的$\pi$介子RITD。
同样，没有可见的晶格间距或$\pi$介子质量依赖。


\begin{figure}[htbp]
\centering
	\centering
	\includegraphics[width=0.45\textwidth]{images/ITD_etas_p456_comparision.pdf}
	\includegraphics[width=0.45\textwidth]{images/ITD_pion_p334_comparision.pdf}
\caption{a12m220、a12m310 与 a15m310 系综在提升动量分别约为2~GeV（上）与1.3~GeV（下）处的$\eta_s$与$\pi$介子RITD。
两种情形下我们都观察到弱的晶格间距与$\pi$介子质量依赖。
}
\label{fig:RITD-ensembles}
\end{figure}


要提取胶子PDF，我们遵循第~\ref{sec:cal-details}节中式~\ref{eq:ME_unpol}至式~\ref{evolution}的步骤：先得到EITD，再用式~\ref{conversion}提取$g(x)$。
为得到EITD，需要RITD $\mathscr{M}(\nu,z^2)$是$\nu$的连续函数，以便计算式~\ref{evolution}中$x\in[0,1]$的积分。
为此，我们对RITD作``$z$ 展开''拟合~\cite{Boyd:1994tt,Bourrely:2008za}%
\footnote{注意，``$z$ 展开''中的$z$与本文其他地方使用的 Wilson 线长度$z$无关。}，
采用文献~\cite{Joo:2019bzr}的形式：
%
\begin{equation}
\mathscr{M}(\nu, z^2,M_{\pi}) = \sum_{k=0}^{k_\text{max}}\lambda_k\tau^k,
\label{zexpansion}
\end{equation}
%
其中$\tau=\frac{\sqrt{\nu_\text{cut}+\nu}-\sqrt{\nu_\text{cut}}}{\sqrt{\nu_\text{cut}+\nu}+\sqrt{\nu_\text{cut}}}$。
随后把拟合的$\mathscr{M}(\nu, z^2)$代入式~\ref{evolution}的积分。
演化函数中$\mathscr{M}(u\nu, z^2)$项的$z$依赖来自一圈匹配项，相对树图项是高阶修正；
因此可忽略$\mathscr{M}(\nu, z^2)$中的$z$依赖。
我们取$\nu_\text{cut}=1$为最佳值（同文献~\cite{Joo:2019bzr}），同时在$[0.5,2]$范围内变动$\nu_\text{cut}$，结果一致。
我们固定$\lambda_0 = 1$以满足式~\ref{eq:RITD}中RITD $\mathscr{M}(\nu, z^2)$的归一化。
展开阶取$k_\text{max}=3$，因为对每个系综使用4项$z$展开即可在$\chi^2/\text{dof}<1$下拟合全部$P_z\in [1,5]\times 2\pi/L$（a12m220 系综为$P_z\in [1,7]\times 2\pi/L$）
且$z$达0.6~fm的数据点。
``$z$ 展开''在RITD上的重构带示于图~\ref{fig:ITD-comparison}上图；
它们对所有系综的RITD数据点都描述良好。


\begin{figure}[htbp]
\centering
	\centering
	\includegraphics[width=0.45\textwidth]{images/RITD_pion_comparision.pdf}
	\includegraphics[width=0.45\textwidth]{images/EITD_pion_comparision.pdf}
\caption{
上图为带``$z$ 展开''拟合重构带的RITD $\mathscr{M}$，下图为带拟合重构带的EITD $G$；计算所用系综为晶格间距$a\approx 0.12$~fm、$\pi$介子质量$M_\pi\approx\{220,310,690\}$~MeV，以及$a\approx 0.15$~fm、$M_\pi\approx 310$~MeV；注意此处$a\approx 0.12$~fm、$M_\pi\approx 690$~MeV 的结果来自 a12m220 系综。
}
\label{fig:ITD-comparison}
\end{figure}


得到连续$\nu$的拟合RITD后，经由式~\ref{evolution}获得EITD。
本研究所有系综的RITD $\mathscr{M}$与EITD $G$随$\nu$的变化示于图~\ref{fig:ITD-comparison}。
在某些$\nu$值处，同一$\nu$对应多组$z$与$P_z$组合；
因此同一颜色与符号的点会来自同一晶格间距与$\pi$介子质量并在同一$\nu$处重叠。
要通过式~\ref{conversion}匹配到光锥胶子PDF，EITD $G(\nu,\mu)$应不含$z^2$依赖；
然而由式~\ref{evolution}所得EITD因忽略匹配中胶子内嵌于夸克的贡献及高阶项而带有$z^2$依赖。
EITD还依赖于晶格间距$a$与$\pi$介子质量$M_\pi$。
回顾RITD对晶格间距$a$和$\pi$介子质量$M_\pi$的依赖较弱，
可见$a$与$M_\pi$依赖对EITD的影响也不大：
如图~\ref{fig:ITD-comparison}第二行所示，不同$a$、$M_\pi$的EITD结果大体彼此一致。
此外我们在图~\ref{fig:ITD-comparison}中还观察到RITD与EITD对$z^2$的依赖很弱。

现在可以用式~\ref{conversion}从EITD提取胶子PDF $g(x,\mu^2)$。
%
我们假设光锥PDF取 JAM~\cite{Barry:2018ort,Cao:2021aci}也采用的函数形式来拟合EITD：
%
\begin{align}
f_g(x,\mu) = \frac{xg(x, \mu)}{\langle x \rangle_g(\mu)} = \frac{x^A(1-x)^C}{B(A+1,C+1)},
\label{functional}
\end{align}
%
适用于$x\in[0,1]$，区间外为零。
贝塔函数$B(A+1,C+1)=\int_0^1 dx\, x^A(1-x)^C$用于把面积归一到一。
然后施加匹配公式，经式~\ref{conversion}从该函数形式PDF得到EITD ${G}$。
我们把参数化给出的EITD $G(\nu,\mu)$与格点计算得到的EITD $G(\nu,z^2,\mu,a,M_\pi)$进行拟合比较。


\begin{figure}[htbp]
\centering
\centering
\includegraphics[width=0.45\textwidth]{images/xg-x0-1_pion_a12m220_par.pdf}
\caption{
晶格间距$a\approx 0.12$~fm、$\pi$介子质量$M_\pi\approx 220$~MeV、$\mu^2=4\text{ GeV}^2$处的$xg(x, \mu)/\langle x \rangle_g$随$x$的分布（底部），并附有按式~\ref{functional}及其后段落所述的1、2与3参数拟合得到的$z_\text{max}\approx 0.6$~fm拟合带。
}
\label{fig:fitform-comparison}
\end{figure}


我们研究了PDF全局分析与部分格点计算中对$f_g(x,\mu)$常用的不同参数化形式所引入的系统不确定度。
第一种是式~\ref{functional}的双参数形式。
其次考虑 xFitter 分析~\cite{Novikov:2020snp}（亦见文献~\cite{Aurenche:1989sx,Sutton:1991ay}）使用的单参数形式$N_1 (1-x)^C$，它等价于式~\ref{functional}取$A=0$。
第三种是三参数形式$N_3 x^A(1-x)^C(1+D\sqrt{x})$。
通过对1、2、3参数PDF形式施加方案转换（式~\ref{conversion}），我们把三种不同形式拟合到$z_\text{max}\approx 0.6$~fm格点数据的EITD。
这里聚焦最轻$\pi$介子质量$M_{\pi}\approx 220$~MeV、晶格间距$a\approx 0.12$~fm的结果。
拟合的$\chi^2/\text{dof}$依次由$1.47(72)$、$1.08(68)$降至$1.04(41)$，表明双参数与三参数拟合的质量略佳。
如图~\ref{fig:fitform-comparison}所示，单参数拟合与双参数拟合的$f_g(x,\mu)$拟合带在$x<0.4$区差异很大，而双参数与三参数拟合之间的差异小得多。
因此我们得出结论：此格点数据上的单参数拟合不太可靠，而拟合结果在双参数与三参数拟合处收敛。
同样的结论对所有其他系综与$\pi$介子质量均成立。
因此最终结果采用式~\ref{functional}定义的双参数形式（与 JAM 相同的参数化）是非常合理的。


系统不确定度的另一来源是：基于（受全局拟合启发的）假设——$\pi$介子$q_S(x)$小于胶子PDF——而在式~\ref{matching-eq}中忽略夸克项的贡献。
目前尚无来自格点模拟的$q_S(x)$结果，因为此前只做过$\pi$介子的价分布。
于是我们用 xFitter~\cite{Novikov:2020snp}的 NLO $\pi$介子夸克PDF来估计省略$q_S(x)$贡献带来的系统误差。
以$a\approx 0.12$~fm、$\pi$介子质量$M_\pi\approx 220$~MeV格点为例，利用这些夸克PDF得到含胶子内嵌于夸克$R_{gq}$项修正后的RITD与EITD，重复自式~\ref{zexpansion}起的同一流程并用式~\ref{functional}拟合EITD。
在图~\ref{fig:xgx}右侧，我们展示同时包含胶子内嵌于胶子（gg）与胶子内嵌于夸克（gq）贡献的$xg(x, \mu)/\langle x \rangle_g$平均值（蓝色实线），并与只用gg贡献的 a12m220 结果（绿色实线）比较。
对$x<0.9$，含gg贡献的平均值存在5\%到$10\%$的差异，这表明在小-$x$区当前统计误差的水平下，$\mu^2=4\text{ GeV}^2$处的胶子内嵌于夸克的贡献相对较小；
而在$x>0.9$区，该贡献变得更显著，但仍小于统计误差。
一旦将来的研究在大-$x$区具有充分减小的统计不确定度，就需要把夸克贡献包括进来。

综合上述对拟合形式选择与夸克项贡献的分析，我们得出结论：相对于当前统计量这些系统误差可以忽略。
因此我们对所有格点系综的最终结果采用$z_\text{max}\approx 0.6$~fm（a15m310 系综为$z_\text{max}\approx 0.75$~fm）的EITD拟合、忽略匹配中的夸克贡献项，并使用式~\ref{conversion}的拟合形式。
各系综重构的$xg(x, \mu)/\langle x \rangle_g$拟合带示于图~\ref{fig:xgx}左图，比较了不同晶格间距与$\pi$介子质量的结果。
在同一晶格间距$a\approx 0.12$~fm下，不同$\pi$介子质量$M_\pi\approx\{220,310,690\}$~MeV的重构拟合带彼此一致，表明胶子PDF对$\pi$介子质量的依赖温和；
类似地，围绕$\pi$介子质量$M_\pi\approx 310$~MeV比较晶格间距依赖时，我们发现0.12~fm格点的拟合PDF在$x>0.1$区略小，但仍在一个标准差之内，表明晶格间距依赖也是温和的。
我们还注意到，在大-$x$区不同系综的带随$x\to 1$的衰减快慢不同；
下面更深入地研究这种衰减行为。

胶子PDF在大-$x$区的衰减行为在理论与全局分析中都被广泛研究。
微扰QCD研究~\cite{Close:1977qx,Sivers:1982wk}与 DSE 计算~\cite{Bednar:2018mtf,Freese:2021zne}表明当$x\to 1$时$g(x,\mu^2) \sim (1-x)^C$，其中$C\approx 3$。
微扰QCD的预言~\cite{Sivers:1982wk}基于这样的观点：由于夸克是大-$x$胶子的源，胶子PDF在大$x$处相对夸克PDF受到压低；
也就是说，$x\to 1$时$g(x,\mu^2)/q_v(x,\mu^2)\to 0$。
早期的实验数据拟合给出$C\approx 2$~\cite{Aurenche:1989sx,Sutton:1991ay}或$C<2$~\cite{Gluck:1991ey,Gluck:1999xe}，而较新的 JAM 合作组全局分析给出$C>3$~\cite{Barry:2018ort,Cao:2021aci}，xFitter 合作组则发现$C\approx 3$~\cite{Novikov:2020snp}。
在晶格间距$a\approx 0.12$~fm下，我们的拟合参数$C$对$M_\pi\approx\{690,310,220\}$~MeV分别为$3.6(1.5)$、$3.3(2.0)$、$4.7(2.8)$。
这些$C$的结果彼此一致，且随着$\pi$介子质量趋近物理$\pi$介子质量呈轻微上升趋势。
对于晶格间距$a\approx\{0.15,0.12\}$~fm，在$M_\pi\approx 310$~MeV下$C=\{2.2(1.5),3.3(2.0)\}$，提示$C$将向连续极限增大。
我们也考察了胶子内嵌于夸克的贡献对$C$值的影响，其造成约$0.1$的差异，予以忽略。
鉴于$\pi$介子质量外推与晶格间距外推似乎都显示$C$增大，由这项格点QCD研究得出$C>3$的结论看来是合理的。

我们在图~\ref{fig:xgx}右图中把重构的胶子PDF与全局拟合结果进行比较。
图中给出$a\approx 0.12$~fm、$M_\pi\approx 220$~MeV格点的$xg(x, \mu)/\langle x \rangle_g$重构拟合带，以及$\mu^2=4\text{ GeV}^2$处 xFitter~\cite{Novikov:2020snp}与 JAM~\cite{Barry:2018ort,Cao:2021aci}的 NLO $\pi$介子胶子PDF。
JAM 带看起来比预期偏宽，因为我们用$xg(x, \mu)$除以$\langle x \rangle_g$的平均值来重构它；
正确估计误差所需的相关联数值无法获得。
注意 xFitter 采用式~\ref{functional}取$A=0$的拟合形式。
我们的拟合$\pi$介子胶子PDF在$x>0.2$与 JAM 一致、在$x>0.5$与 xFitter 在一个标准差内一致。
从这一比较可见，我们结果的误差大小与全局分析相近，能够在实验数据之外提供来自理论计算的约束。


\begin{figure*}[htbp]
\centering
	\centering
	\includegraphics[width=0.45\textwidth]{images/xg-x0-1_bandcomp_aM_pion.pdf}
	\includegraphics[width=0.45\textwidth]{images/xg-x0-1_bandcomp_JX_pion.pdf}
\caption{
由格点数据拟合得到的$\pi$介子胶子PDF $xg(x, \mu)/\langle x \rangle_g$随$x$的分布：左图为晶格间距$a\approx\{0.12,0.15\}$~fm、$\pi$介子质量$M_\pi\approx\{220,310,690\}$~MeV各系综的结果；右图为$a\approx 0.12$~fm、$M_\pi\approx 220$~MeV格点的结果，并与$\mu=2$~GeV、$\overline{\text{MS}}$方案下 xFitter 与 JAM 的 NLO $\pi$介子胶子PDF比较。
由于缺少组成分量间可用的相关不确定度，所示 JAM 误差被高估。
我们的PDF结果在$x>0.2$与 JAM~\cite{Barry:2018ort,Cao:2021aci}、$x>0.5$与 xFitter~\cite{Novikov:2020snp}一致。
}
\label{fig:xgx}
\end{figure*}

